\section{From harmonic differences to local metric extension}\label{sec:bootstrap}

Lemma~\ref{local:density} bounds the transfer density uniformly on the
shrinking spheres.  Exact reproduction of affine functions now turns that
bound into a two-power estimate for the difference of paired harmonic
functions.  The harmonic Hessians convert derivatives of these differences
into metric derivatives.  Finite-order interpolation will then improve
the temporary metric bound and allow the radii to shrink with the spatial
dilation held fixed.

\subsection{From the functional to harmonic jets}

Return to the finite harmonic family in Theorem~\ref{local:extension}.
Its values and first derivatives at $0$ agree between the two metrics,
because the functions agree on the one-sided open set and are smooth.
After dilation define
\begin{equation}\label{local:renormalized-pairs}
 v_{j,\lambda}^a(y)=\lambda^{-2}
 \bigl(u_j^a(\lambda y)-u_1^a(0)
                     -\lambda\,du_1^a(0)\cdot y\bigr).
\end{equation}
The subtracted functions are harmonic because the coordinates are
harmonic.  These new pairs still agree on the known side and are
transferred by $T_y$.  For each fixed integer $m$ they have uniform
$C^m$ bounds on every required fixed bounded coordinate neighborhood.
For derivative orders $j\geq2$, this follows from the factor
$\lambda^{j-2}$; orders zero and one follow from the second-order
Taylor remainder.  Moreover,
\[
 \partial^2v_{1,\lambda}^a(y)\longrightarrow\partial^2u_1^a(0)
\]
smoothly on fixed sets.  Thus the Hessian spanning condition has a
uniform rank bound after sufficiently small dilation.  Write
$V^a=v_{2,\lambda}^a-v_{1,\lambda}^a$.

\begin{lemma}[Two-power moment bound]\label{local:moment}
Under the coefficient condition \eqref{balls:bootstrap-condition},
the paired differences satisfy
\begin{equation}\label{local:moment-bound}
 \left\|\frac{\chi V^a}{r_\delta^2}\right\|_{L^2(Y)}\leq C
\end{equation}
with a constant independent of $\delta$ and sufficiently small
$\lambda$.
\end{lemma}

\begin{proof}
Fix a center $y$ and subtract from $v_{1,\lambda}^a$ its affine
coordinate Taylor polynomial at $y$.  Exact constant and coordinate
reproduction in \eqref{balls:reproduction} cancels this polynomial in
$T_yv_{1,\lambda}^a-v_{1,\lambda}^a(y)=V^a(y)$.  The remainder on
$\partial B_\gamma(y,r)$ is bounded by $Cr^2$, since coordinate
distance and $\gamma$-distance are uniformly comparable and the
functions have uniform second derivatives.  Hence
\[
 |V^a(y)|\leq C r_\delta(y)^2
                         \|\psi(y,\cdot)\|_{L^2(S_y)}.
\]
Lemma~\ref{local:density} proves \eqref{local:moment-bound}.  The
coordinate volume measure and $dV_\gamma$ are uniformly comparable.
No derivative in $y$ of the continued density is used here.
\end{proof}

We next express metric derivatives in terms of derivatives of the paired
differences.  Recovery of the conformal metric
class from the hyperplane of harmonic Hessians is also used in
\citep[Proposition~4.1 and Remark~4.2]{LLS20}.
All Hessians in this step are
ordinary coordinate Hessians, not covariant Hessians.  Because the
coordinates are harmonic, the equations for $v_{j,\lambda}^a$ have
no first-order term in these coordinates:
\[
 (g_{k,\lambda}^{-1})^{ij}\partial_{ij}v_{k,\lambda}^a=0
 \qquad(k=1,2).
\]
Multiplication by the scalar factors in
\eqref{balls:normalization} therefore gives
\begin{equation}\label{local:metric-system}
 h^{ij}\partial_{ij}v_{1,\lambda}^a
       =-Q^{ij}\partial_{ij}V^a,
 \qquad \operatorname{tr}(\gamma h)=0.
\end{equation}

Harmonicity puts each $\partial^2v_{1,\lambda}^a(y)$ in the hyperplane
annihilated by $\gamma^{-1}(y)$.  The uniform rank bound established
above shows that these matrices span this hyperplane for small $\lambda$.  The extra equation
$\operatorname{tr}(\gamma h)=0$ completes an injective system on
symmetric contravariant matrices: an element annihilating all those
Hessians is a multiple of $\gamma^{-1}$, and its $\gamma$-trace is
$n$ times that multiple.  The associated finite matrix has a smooth
left inverse with uniform derivatives on $Y$.  For example, in fixed
coordinate bases use $(A^*A)^{-1}A^*$, where $A$ is this injective
matrix.  Its smallest singular value stays bounded below by the
spanning condition and smooth convergence under dilation.
Differentiating \eqref{local:metric-system} at most eight times gives
\begin{equation}\label{local:metric-jet-bound}
 H_8(y)\leq C\sum_a\sum_{|\alpha|\leq10}
                              |\partial^\alpha V^a(y)|.
\end{equation}
The derivatives of $Q$, $\gamma$, and of the finite harmonic family
are uniformly bounded.  In particular, this inversion introduces no
factor involving $r_\delta^{-1}$.

The metric estimate requires derivatives of the paired differences
through order ten, whereas Lemma~\ref{local:moment} controls only their
weighted $L^2$ norm.  The following finite-order interpolation estimate
bridges these two statements using the already fixed smooth bounds.

\begin{lemma}[Finite-order interpolation]\label{local:interpolation}
Let $V$ be smooth on a neighborhood of $\overline{B(y,s)}\subset
\mathbb R^n$, with $0<s\leq1$ and $\|V\|_{C^m}\leq A_m$ there.
For every multi-index $\alpha$ of length $j<m$,
\begin{equation}\label{local:interpolation-estimate}
 |\partial^\alpha V(y)|\leq C_{m,n}
 \bigl(s^{-j-n/2}\|V\|_{L^2(B(y,s))}
                          +A_m s^{m-j}\bigr).
\end{equation}
\end{lemma}

\begin{proof}
Expand $V(y+sz)$ at $z=0$ through degree $m-1$ on the unit ball.
The remainder has absolute value at most $C_{m,n}A_m s^m$.
The $L^2$ norm of the Taylor polynomial is therefore at most
$s^{-n/2}\|V\|_{L^2(B(y,s))}+C_{m,n}A_m s^m$.
On the finite-dimensional space of polynomials of degree at most
$m-1$, every coefficient is bounded by a constant times the $L^2$
norm on the unit ball.  Its $z^\alpha$ coefficient is
$s^j\partial^\alpha V(y)/\alpha!$, giving the claim.
\end{proof}

\subsection{A strict improvement of the metric bound}

Set
\begin{equation}\label{local:exponents}
 D=10+n/2,\qquad \eta=\frac1{2D},\qquad
 p>2D,\qquad m\geq\lceil10+3D\rceil.
\end{equation}
We also take the integer $p$ large enough for
Proposition~\ref{angular:coercivity}.  Thus $\eta>1/p$, and $m$ is
one fixed finite smoothness order.  The constants may depend on these
choices, which precede the final spatial dilation and all shrinkage.

\begin{lemma}[Superlinear improvement]\label{local:improvement}
After the choices in \eqref{local:exponents}, there are $C>0$ and
$\rho_0>0$, independent of $\delta$ and sufficiently small $\lambda$,
such that the coefficient condition
\eqref{balls:bootstrap-condition} implies
\begin{equation}\label{local:superlinear}
 H_8(y)\leq C r_\delta(y)^{3/2}
 \quad\text{for }y\in K\text{ with }r_\delta(y)<\rho_0.
\end{equation}
\end{lemma}

\begin{proof}
Put $r=r_\delta(y)$ and $s=r^\eta$.  By
Lemma~\ref{balls:profiles}, $r_\delta^{1/p}$ has a fixed Lipschitz
constant.  Since
\[
 \frac{s}{r^{1/p}}=r^{\eta-1/p}\longrightarrow0,
\]
a sufficiently small threshold $\rho_0$ ensures
\[
 \tfrac12 r^{1/p}\leq r_\delta(z)^{1/p}
                         \leq\tfrac32 r^{1/p}
             \quad(z\in B(y,s)).
\]
The same fixed threshold ensures that this ball lies where $\chi=1$,
because $K$ has a fixed buffer inside that set.  The upper comparison
and Lemma~\ref{local:moment} imply
\[
 \|V^a\|_{L^2(B(y,s))}\leq C r^2.
\]
Apply Lemma~\ref{local:interpolation} and the fixed uniform $C^m$
bounds of the renormalized functions.  For $j\leq10$ it gives
\[
 |\partial^\alpha V^a(y)|
   \leq C_m\bigl(r^{2-\eta(j+n/2)}+r^{\eta(m-j)}\bigr)
   \leq C_m r^{3/2}.
\]
Both exponents are at least $3/2$ by \eqref{local:exponents}.
Equation~\eqref{local:metric-jet-bound} completes the proof.
\end{proof}

The gain in \eqref{local:superlinear} is stronger than the temporary
bound $H_8\leq\varepsilon r_\delta$.  Its usefulness is that a single
choice of the spatial dilation makes the improvement strict at every
radius, including those above the small threshold.

\subsection{Local extension of the isometry}

\begin{proof}[Proof of Theorem~\ref{local:extension}]
Make the linear normalization and spatial dilation of
Section~\ref{sec:balls}.  All harmonic functions and their finite
spanning family are chosen before the dilation.  We now specify the
remaining choices and hold each one fixed before making the next.

First choose the fixed conformal geometry, cylinder, depth function,
and cutoff, keeping uniform background smooth bounds for
$0<\lambda\leq\lambda_0$.  Choose $p,\eta,m$ as in
\eqref{local:exponents}, with $p$ large enough for the geometric and
angular estimates.  Choose $\varepsilon$ for those estimates and then
$R_*$ so that the balls lie in normal neighborhoods and
$r_\delta\beta_\delta^2\leq\varepsilon$ for every $\delta$.
These choices fix every separation range needed to start the
continuation and to bound the cutoff terms.  Restrict $\lambda_0$
further so that Proposition~\ref{cross:initial} applies to these fixed
ranges, with bounds uniform for $0<\lambda\leq\lambda_0$.

The constants $C,\rho_0$ of Lemma~\ref{local:improvement} are now
fixed.  Choose
\begin{equation}\label{local:rho-choice}
 0<\rho<\rho_0,\qquad C\rho^{1/2}\leq\varepsilon/2.
\end{equation}
At points with $r_\delta<\rho$, that lemma improves the coefficient
bound to $H_8\leq\varepsilon r_\delta/2$ whenever the temporary bound
holds.  At points with $r_\delta\geq\rho$, the same improvement will
hold if
\begin{equation}\label{local:one-dilation}
 \sup_Y H_8\leq\varepsilon\rho/2.
\end{equation}
Since $h\to0$ smoothly under dilation, choose one
$\lambda\leq\lambda_0$ satisfying \eqref{local:one-dilation} and the
additional initial bounds described next.  This $\lambda$ is kept
fixed for the rest of the proof.

There is a uniform reason the coefficient bound holds off $K$ for
every $\delta$.  Outside $K$, either $y_n<-1/2$, where $h=0$ after
the fixed dilation is sufficiently small, or $t>2t_b$, where
\[
 r_\delta\geq R_*(2t_b)^p>0.
\]
Smooth smallness of $h$ therefore gives
$H_8\leq\varepsilon r_\delta/2$ off $K$ for every $\delta$ by
reducing the same dilation once.  At $\delta=1$, the radius has a
positive minimum on all of $Y$, so we can also ensure the initial
coefficient bound there.  These conditions involve only fixed
positive constants, and are compatible with
\eqref{local:one-dilation}.

For this fixed metric pair on the dilated neighborhood, set
\[
 \mathcal A=\{\delta\in(0,1]:
        H_8(y)\leq\varepsilon r_\delta(y)\text{ for every }y\in K\}.
\]
It contains $1$.  At any $\delta\in\mathcal A$, the bound already
established off $K$ gives the full hypothesis
\eqref{balls:bootstrap-condition}.  Proposition~\ref{balls:continuation}
and the uniform density estimate of Lemma~\ref{local:density} then apply.
The moment bound, metric system, and interpolation argument give
Lemma~\ref{local:improvement}.
The small-radius improvement \eqref{local:rho-choice} and the
large-radius bound \eqref{local:one-dilation} give
\[
 H_8\leq\tfrac12\varepsilon r_\delta\quad\text{on }K.
\]
For every fixed positive $\delta$, $r_\delta$ has a positive minimum
on the compact set $K$.  Thus
$\delta\mapsto\max_K(H_8/r_\delta)$ is continuous on $(0,1]$.
The set $\mathcal A$ is relatively closed, and the strict improvement
makes it relatively open.  Connectedness implies
$\mathcal A=(0,1]$.

Finally let $\delta\downarrow0$, with the spatial dilation, harmonic
family, cutoff and all structural parameters fixed.  On the fixed
open neighborhood of $0$ where $t<0$, the radii tend to zero.  Hence
$h=0$ there.  By \eqref{balls:normalization}, $g_{2,\lambda}$ is
conformal to $g_{1,\lambda}$ on that neighborhood.  Write
$g_{2,\lambda}=e^{2\omega}g_{1,\lambda}$.  The conformal change formula
for the Laplace--Beltrami operator is
\[
 \Delta_{e^{2\omega}g}u
  =e^{-2\omega}\bigl(\Delta_g u
                 +(n-2)\langle d\omega,du\rangle_g\bigr).
\]
Apply it to every common harmonic coordinate.  Since their
differentials form a basis and $n-2\ne0$, it forces $d\omega=0$.
Since $g_{1,\lambda}(0)=g_{2,\lambda}(0)=I$, this constant is zero.
This is the local harmonic-morphism rigidity mechanism of
\citep[Section~8]{Fuglede1978}, in the restricted-family form used in
\citep[Proposition~4.1]{LLS20}.
The metrics therefore agree on a connected neighborhood of $0$;
undoing the dilation proves the theorem.
\end{proof}

For clarity, the two limits in this proof have different roles:
spatial dilation is used once to obtain a fixed small metric
difference, and the subsequent shrinking family improves the bound
for that same difference.  The parameter order is
\begin{equation}\label{local:parameter-order}
 \begin{gathered}
 \text{fixed geometry and harmonic family}
   \;\longrightarrow\;(p,\eta,m)
   \;\longrightarrow\;(\varepsilon,R_*)\\
 \longrightarrow\;\rho\;\longrightarrow\;
 \lambda\text{ fixed}\;\longrightarrow\;\delta\downarrow0.
 \end{gathered}
\end{equation}
All estimates use finitely many derivatives selected before either
limiting step.  Theorem~\ref{local:extension} is the local input that
will prevent a maximal matched region from having an interior
frontier.
